\documentclass[11pt]{article}

\usepackage[margin=1in]{geometry}
\usepackage[T1]{fontenc}
\usepackage{lmodern}
\usepackage{microtype}
\usepackage{amsmath,amssymb}
\usepackage[dvipsnames]{xcolor}
\usepackage{enumitem}
\usepackage[most]{tcolorbox}
\usepackage[hidelinks]{hyperref}

\hypersetup{
  pdftitle={Homework 1 Solutions, MATH 327, Fall 2026},
  pdfauthor={Manish Patnaik},
  pdfsubject={MATH 327 homework solutions}
}

\definecolor{courseblue}{HTML}{1F4E79}
\definecolor{lightblue}{HTML}{EAF2F8}
\definecolor{boxborder}{HTML}{B8CBDC}
\definecolor{solutiongreen}{HTML}{2E6B4F}
\definecolor{lightsolution}{HTML}{F2F8F4}
\definecolor{darkgray}{HTML}{333333}

\setlength{\parindent}{0pt}
\setlength{\parskip}{0.55em}
\setlist[enumerate]{leftmargin=2em, itemsep=0.3em, topsep=0.3em}

\newtcolorbox{exercise}[1]{
  enhanced, breakable,
  colback=white,
  colframe=boxborder,
  colbacktitle=lightblue,
  coltitle=courseblue,
  fonttitle=\bfseries,
  title={Exercise #1},
  boxrule=0.6pt,
  arc=1mm,
  left=2.5mm, right=2.5mm, top=1.5mm, bottom=1.5mm,
  before skip=0.8em, after skip=0.35em
}

\newtcolorbox{solution}{
  enhanced, breakable,
  colback=lightsolution,
  colframe=solutiongreen!55,
  colbacktitle=solutiongreen!12,
  coltitle=solutiongreen,
  fonttitle=\bfseries,
  title={Solution},
  boxrule=0.6pt,
  arc=1mm,
  left=2.5mm, right=2.5mm, top=1.5mm, bottom=1.5mm,
  before skip=0.25em, after skip=0.9em
}

\newcommand{\chapterheading}[1]{%
  \par\vspace{0.5em}
  {\color{courseblue}\Large\bfseries #1}\par\vspace{0.35em}
  {\color{boxborder}\hrule height 0.6pt}\vspace{0.35em}
}

\newcommand{\sectionheading}[1]{%
  \vspace{0.25em}
  {\color{darkgray}\large\bfseries #1}\par
}

\newcommand{\R}{\mathbb{R}}
\newcommand{\Z}{\mathbb{Z}}
\newcommand{\GL}{\operatorname{GL}}
\newcommand{\im}{\operatorname{im}}

\begin{document}

\begin{center}
  \colorbox{lightblue}{%
    \begin{minipage}{0.92\textwidth}
      \vspace{0.6em}
      \begin{tabular*}{\textwidth}{@{\extracolsep{\fill}}lr}
        {\color{courseblue}\bfseries\LARGE Homework 1 Solutions} &
        {\color{courseblue}\bfseries\LARGE MATH 327}\\[0.4em]
        \multicolumn{2}{l}{\color{darkgray}\bfseries Fall 2026}
      \end{tabular*}
      \vspace{0.5em}
    \end{minipage}}
\end{center}

\vspace{0.5em}

Unless otherwise indicated, the following exercises are from Michael Artin's
\emph{Algebra}, second edition.  We write products of permutations using the
right-to-left convention, and put $s_i=(i\ i+1)$.

\chapterheading{Chapter 1: Matrices}

\sectionheading{Section 1: The Basic Operations}

\begin{exercise}{1.13}
A square matrix $A$ is \emph{nilpotent} if $A^k=0$ for some $k>0$.
Prove that if $A$ is nilpotent, then $I+A$ is invertible. Do this by
finding the inverse.
\end{exercise}

\begin{solution}

The key intuition is to apply the geometric series formula formally to matrices:
\[
  (I+A)^{-1}=I-A+A^2-A^3+\cdots.
\]
In general, the infinite series on the right-hand side need not converge.
However, since $A$ is nilpotent, we can choose $k>0$ such that $A^k=0$.
All powers of degree at least $k$ then vanish, suggesting the finite sum
\[
  S=I-A+A^2-A^3+\cdots+(-1)^{k-1}A^{k-1}.
\]
This is a finite geometric sum.  Multiplying on either side gives
\[
 (I+A)S=S(I+A)=I+(-1)^{k-1}A^k=I.
\]
Thus $I+A$ is invertible and
\[
  \boxed{(I+A)^{-1}=I-A+A^2-\cdots+(-1)^{k-1}A^{k-1}}.
\]
\end{solution}

\sectionheading{Miscellaneous Problems}

\begin{exercise}{M.10}
Let $A,B$ be $m\times n$ and $n\times m$ matrices. Prove that
$I_m-AB$ is invertible if and only if $I_n-BA$ is invertible.
\end{exercise}

\begin{solution}
Suppose first that $I_m-AB$ is invertible. Using the same intuition as in
the previous exercise, its inverse should formally be
\[
  I_m+AB+(AB)^2+(AB)^3+\cdots.
\]
Multiplying this series by $B$ on the left and $A$ on the right, still
formally, gives
\[
  BA+(BA)^2+(BA)^3+\cdots.
\]
Adding $I_n$ gives the formal geometric series for $(I_n-BA)^{-1}$.
These series need not converge, but they suggest a formula that we can
verify directly.

We first \textit{define}
$X=(I_m-AB)^{-1}$ and then consider
\[
  Y=I_n+BXA.
\]
We claim this is the desired inverse of $I_n-BA$.  Indeed,
\begin{align*}
 (I_n-BA)Y
 &=I_n-BA+BXA-BABXA\\
 &=I_n-BA+B(I_m-AB)XA=I_n,
\end{align*}
because $(I_m-AB)X=I_m$.  Similarly,
\begin{align*}
 Y(I_n-BA)
 &=I_n-BA+BXA-BXABA\\
 &=I_n-BA+BX(I_m-AB)A=I_n,
\end{align*}
because $X(I_m-AB)=I_m$.  Therefore
\[
  (I_n-BA)^{-1}=I_n+B(I_m-AB)^{-1}A.
\]
The converse follows by interchanging $A$ and $B$ (and $m$ and $n$).
\end{solution}

\chapterheading{Chapter 2: Groups}

\sectionheading{Supplementary Problems: Permutations, Cycles, and Length}

\begin{exercise}{P.1}
Let $p=[4,1,6,2,5,3]\in S_6$.
\begin{enumerate}[label=\textbf{(\alph*)}]
  \item Write $p$ as a product of disjoint cycles.
  \item List all inversions of $p$.
  \item Compute $\operatorname{inv}(p)$ and hence determine $\ell(p)$.
  \item Express $p$ as a product of exactly $\ell(p)$ simple transpositions.
\end{enumerate}
\end{exercise}

\begin{solution}
\begin{enumerate}[label=\textbf{(\alph*)}]
  \item Following the images of the entries gives
  \[
    1\mapsto4\mapsto2\mapsto1,
    \qquad 3\mapsto6\mapsto3,
    \qquad 5\mapsto5.
  \]
  Hence $p=(1\ 4\ 2)(3\ 6)$, with the fixed point $5$ omitted.

  \item An inversion is a pair $(i,j)$ with $i<j$ and $p(i)>p(j)$.
  Reading from the one-line notation, the inversions are
  \[
  (1,2),(1,4),(1,6),
  (3,4),(3,5),(3,6),(5,6).
  \]

  \item There are seven inversions, so
  \[
    \operatorname{inv}(p)=7
    \qquad\text{and hence}\qquad
    \ell(p)=7.
  \]

  \item Starting with the identity in one-line notation and successively
  swapping adjacent positions gives
  \begin{align*}
  [1,2,3,4,5,6]
  &\xrightarrow{s_3}[1,2,4,3,5,6]
  \xrightarrow{s_2}[1,4,2,3,5,6]
  \xrightarrow{s_1}[4,1,2,3,5,6]\\
  &\xrightarrow{s_5}[4,1,2,3,6,5]
  \xrightarrow{s_4}[4,1,2,6,3,5]
  \xrightarrow{s_3}[4,1,6,2,3,5]\\
  &\xrightarrow{s_5}[4,1,6,2,5,3].
  \end{align*}
  Thus
  \[
    \boxed{p=s_3s_2s_1s_5s_4s_3s_5}.
  \]
  This expression has seven factors, so it is reduced.
\end{enumerate}
\end{solution}

\begin{exercise}{P.2}
Let $1\leq a<b\leq n$, and let $\tau=(a\ b)\in S_n$.
\begin{enumerate}[label=\textbf{(\alph*)}]
  \item Write $\tau$ in one-line notation.
  \item Prove that $\operatorname{inv}(\tau)=2(b-a)-1$.
  \item Deduce that $\ell(\tau)=2(b-a)-1$.
  \item Explain why two permutations with the same cycle type need not have
        the same Coxeter length.
\end{enumerate}
\end{exercise}

\begin{solution}
\begin{enumerate}[label=\textbf{(\alph*)}]
  \item The transposition fixes every entry except $a$ and $b$, so its
  one-line notation is
  \[
  [1,\ldots,a-1,\,b,\,a+1,\ldots,b-1,\,a,\,b+1,\ldots,n].
  \]

  \item The inversions involving position $a$ are
  \[
    (a,a+1),(a,a+2),\ldots,(a,b),
  \]
  giving $b-a$ inversions.  The remaining inversions are
  \[
    (a+1,b),(a+2,b),\ldots,(b-1,b),
  \]
  giving $b-a-1$ more.  All other pairs retain their natural order.
  Therefore
  \[
    \operatorname{inv}(\tau)=(b-a)+(b-a-1)=2(b-a)-1.
  \]

  \item Since Coxeter length in $S_n$ equals the number of inversions,
  \[
    \boxed{\ell(\tau)=2(b-a)-1}.
  \]

  \item All transpositions have the same cycle type, namely one $2$-cycle
  and $n-2$ fixed points, but their lengths depend on the distance between
  the exchanged positions.  For example, when $n\geq3$,
  \[
    \ell((1\ 2))=1,
    \qquad
    \ell((1\ n))=2n-3.
  \]
  Thus cycle type does not determine Coxeter length.
\end{enumerate}
\end{solution}

\sectionheading{Section 3: Subgroups of the Additive Group of Integers}

\begin{exercise}{3.3}
\begin{enumerate}[label=\textbf{(\alph*)}]
  \item Define the greatest common divisor of $\{a_1,\ldots,a_n\}$.
        Prove that it exists and is an integer combination of the $a_i$.
  \item If the greatest common divisor is $d$, prove that the greatest
        common divisor of $\{a_1/d,\ldots,a_n/d\}$ is $1$.
\end{enumerate}
\end{exercise}

\begin{solution}
Assume the $a_i$ are not all zero.
\begin{enumerate}[label=\textbf{(\alph*)}]
  \item The greatest common divisor is the unique positive integer $d$
  such that
  \begin{enumerate}[label=(\roman*), leftmargin=2.2em]
    \item $d\mid a_i$ for every $i$, and
    \item every common divisor of the $a_i$ divides $d$.
  \end{enumerate}
  To prove existence, consider
  \[
    H=\Z a_1+\cdots+\Z a_n
     =\{r_1a_1+\cdots+r_na_n:r_i\in\Z\}.
  \]
  This is a nonzero subgroup of $\Z$, so $H=d\Z$ for a unique $d>0$.
  Since $d\in H$, there are integers $r_1,\ldots,r_n$ for which
  \[
    d=r_1a_1+\cdots+r_na_n.
  \]
  Also $a_i\in H=d\Z$, so $\Z a_i\subseteq\Z d$ and hence $d\mid a_i$ for every $i$.  If $c$ is any
  common divisor of the $a_i$, then $c$ divides every integer combination
  of them, and therefore $c\mid d$.  Hence this $d$ is the gcd and is an
  integer combination of the $a_i$.

  \item Put $a_i'=a_i/d$.  Dividing the displayed B\'ezout identity by
  $d$ gives
  \[
    1=r_1a_1'+\cdots+r_na_n'.
  \]
  Any common divisor of the $a_i'$ must divide the left-hand side, so the
  only positive common divisor is $1$.  Thus
  $\gcd(a_1/d,\ldots,a_n/d)=1$.
\end{enumerate}
\end{solution}

\sectionheading{Section 4: Cyclic Groups}

\begin{exercise}{4.3}
Let $a,b\in G$. Prove that $ab$ and $ba$ have the same order.
\end{exercise}

\begin{solution}
The elements are conjugate, since
\[
  ba=b(ab)b^{-1}.
\]
For every integer $r$, conjugation gives
\[
  (ba)^r=b(ab)^rb^{-1}.
\]
Consequently $(ba)^r=1$ if and only if $(ab)^r=1$.  Thus the least
positive such $r$ is the same for both elements; if no such $r$ exists,
both have infinite order.  Hence $ab$ and $ba$ have the same order.
\end{solution}

\begin{exercise}{4.7}
Let $x,y\in G$. Assume that $x$, $y$, and $xy$ each have order $2$.
Prove that $H=\{1,x,y,xy\}$ is a subgroup of $G$ of order $4$.
\end{exercise}

\begin{solution}
Because $x^2=y^2=(xy)^2=1$, each of $x,y,xy$ is its own inverse.  Moreover,
\[
  yx=y^{-1}x^{-1}=(xy)^{-1}=xy,
\]
so $x$ and $y$ commute.  It follows that every product of powers of $x$
and $y$ reduces to one of
\[
  1,\quad x,\quad y,\quad xy.
\]
Thus $H=\langle x,y\rangle$ and is a subgroup.  The four displayed
elements are distinct: $x,y,xy$ are not $1$ because each has order $2$;
$x=y$ would imply $xy=x^2=1$; and $xy=x$ or $xy=y$ would imply respectively
$y=1$ or $x=1$.  Hence $|H|=4$.
\end{solution}

\sectionheading{Section 5: Homomorphisms}

\begin{exercise}{5.3}
Let $U$ be the group of invertible upper triangular matrices
$A=\begin{bmatrix}a&b\\0&d\end{bmatrix}$, and define
$\varphi\colon U\to\R^\times$ by $\varphi(A)=a^2$. Prove that $\varphi$
is a homomorphism, and determine its kernel and image.
\end{exercise}

\begin{solution}
For
\[
 A=\begin{bmatrix}a&b\\0&d\end{bmatrix},\qquad
 C=\begin{bmatrix}c&e\\0&f\end{bmatrix}\in U,
\]
the upper-left entry of $AC$ is $ac$.  Therefore
\[
  \varphi(AC)=(ac)^2=a^2c^2=\varphi(A)\varphi(C),
\]
so $\varphi$ is a homomorphism.  Since an upper triangular matrix is
invertible exactly when $a,d\neq0$,
\[
 \ker\varphi
 =\left\{
 \begin{bmatrix}\varepsilon&b\\0&d\end{bmatrix}:
 \varepsilon\in\{1,-1\},\ b\in\R,\ d\in\R^\times
 \right\}.
\]
Every value $a^2$ is positive, so $\im\varphi\subseteq\R_{>0}$.  Conversely,
if $r>0$, then
\[
 \varphi\!\left(\begin{bmatrix}\sqrt r&0\\0&1\end{bmatrix}\right)=r.
\]
Hence $\im\varphi=\R_{>0}$.
\end{solution}

\begin{exercise}{5.6}
Determine the center of $\GL_n(\R)$.
\end{exercise}

\begin{solution}
Every nonzero scalar matrix commutes with every matrix, so
$\{\lambda I_n:\lambda\in\R^\times\}$ lies in the center.

Conversely, let $A=(a_{pq})$ be central.  For $i\neq j$, the elementary
matrix $T_{ij}=I_n+E_{ij}$ is invertible.  From $AT_{ij}=T_{ij}A$ we get
$AE_{ij}=E_{ij}A$.  Entrywise this says
\[
  a_{pi}\,\delta_{jq}=\delta_{pi}\,a_{jq}
  \qquad\text{for all }p,q.
\]
Taking $q=j$ and $p\neq i$ gives $a_{pi}=0$, so column $i$ has no
off-diagonal entries.  Taking $p=i$ and $q\neq j$ gives $a_{jq}=0$, so
row $j$ has no off-diagonal entries.  Finally, taking $(p,q)=(i,j)$ gives
$a_{ii}=a_{jj}$.  As $i\neq j$ were arbitrary, $A$ is diagonal and all
its diagonal entries are equal.  Since $A$ is invertible, this common
entry is nonzero.  Therefore
\[
  \boxed{Z(\GL_n(\R))=\{\lambda I_n:\lambda\in\R^\times\}}.
\]
(For $n=1$, the same conclusion is immediate.)
\end{solution}

\sectionheading{Section 6: Isomorphisms}

\begin{exercise}{6.7}
Let $H\leq G$ and let $g\in G$. Prove that
$gHg^{-1}=\{ghg^{-1}:h\in H\}$ is a subgroup of $G$.
\end{exercise}

\begin{solution}
The set is nonempty because $g1g^{-1}=1$.  If
$x=gh_1g^{-1}$ and $y=gh_2g^{-1}$ belong to $gHg^{-1}$, then
\[
  xy^{-1}
  =(gh_1g^{-1})(gh_2g^{-1})^{-1}
  =g(h_1h_2^{-1})g^{-1}.
\]
Since $h_1h_2^{-1}\in H$, we have $xy^{-1}\in gHg^{-1}$.  The subgroup
test now shows that $gHg^{-1}\leq G$.
\end{solution}

\begin{exercise}{6.11}
Let $a\in G$. Prove that if $\{1,a\}$ is a normal subgroup of $G$, then
$a$ lies in the center of $G$.
\end{exercise}

\begin{solution}
If $a=1$, the conclusion is immediate.  Otherwise, let $g\in G$.  Normality
gives
\[
  gag^{-1}\in\{1,a\}.
\]
Conjugation is injective, so $gag^{-1}\neq1$ because $a\neq1$.
Consequently $gag^{-1}=a$, and hence $ga=ag$.  This holds for every
$g\in G$, so $a\in Z(G)$.
\end{solution}

\vfill
{\small\color{darkgray}Compiled \today.}

\end{document}
